\chapter{Depth Scaling and Covariance SDE Limits}\label{ch:depth}

These notes study what happens when network \emph{depth} becomes unlimited ($D \in \hNinf$). The recurring template is: (i) identify an infinitesimal scaling (step size) per layer so that each layer contributes an infinitesimal perturbation over a hyperfinite depth grid, and (ii) take the standard part to obtain a continuous-time shadow.

Throughout this chapter, unless stated otherwise, weight entries are i.i.d.\ standard Gaussian (internally $\hNcal(0,1)$ in hyperfinite models) and activations are applied coordinatewise. For each internal model with mutually $\hP$-independent layer weights, let $\mathcal{F}_\ell$ denote the internal algebra generated by all hidden states through layer $\ell$ (including every propagated input when there are several). By Transfer, conditioning on $\mathcal{F}_\ell$ makes the next-layer preactivations internally Gaussian at every hyperfinite width, just as in the finite-width identity~(2.26) of Chapter~2. Unlimited-width shadows enter when the resulting internal empirical covariances and their one-step fluctuations are lifted via the standard-part map to a deterministic recursion or an It\^o diffusion on the Loeb probability space $(\Omega, \mathcal{L}(\mathcal{A}), \Loeb)$ \cite{loeb1975conversion,anderson1976nonstandard,keisler1984infinitesimal}. Although this chapter focuses on networks at initialization with i.i.d.\ Gaussian weight entries, Section~6.3 in Chapter~6 develops a one-step training-time calculation in the same proportional depth--width framework. We use $\varphi$ (with natural extension $\ns{\varphi}$) for the activation and write standard data points as $\{x^\alpha\}_{\alpha=1}^m$.

\section{ResNet Scaling Limits}\label{sec:resnet_scaling}

We begin with the hyperfinite Euler approximation underlying the deterministic-depth regime. An ordinary differential equation (ODE) is determined by a standard vector field $b$ and an initial condition $x_0$.

\begin{definition}[ODE Flow]\label{def:ode_flow}
Let $b : \R^n \to \R^n$ and $x_0 \in \R^n$. An ODE solution is a differentiable path $(X_t)_{t \in [0,T]}$ satisfying
\begin{equation}\label{eq:ode_flow}
\partial_t X_t = b(X_t), \qquad X_0 = x_0.
\end{equation}
\end{definition}

For an unlimited depth $D \in \hNinf$ and positive infinitesimal step size $\eta = T/D \approx 0$ (or more generally any positive $\eta \approx 0$) on the hyperfinite layer-time grid $\mathbb{T}_D = \{0, \eta, 2\eta, \dots, D\eta\}$, the simplest numerical scheme is the internal Euler recursion driven by the natural extension $\ns{b} : \hR^n \to \hR^n$:
\begin{equation}\label{eq:euler_scheme}
Y_{k+1} = Y_k + \eta\,\ns{b}(Y_k), \qquad Y_0 = x_0.
\end{equation}
It is helpful to compare the hyperfinite sequence to its piecewise-constant interpolation:
\begin{equation}\label{eq:euler_interp}
Y_t^{(\eta)} := Y_{\lfloor t/\eta \rfloor}, \qquad t \in [0, T].
\end{equation}

\begin{theorem}[Euler Convergence]\label{thm:euler_conv}
Assume $b$ is globally Lipschitz on $\R^n$:
\begin{equation}\label{eq:b_lipschitz}
\norm{b(x) - b(y)} \le C_1 \norm{x - y} \qquad \forall x, y \in \R^n.
\end{equation}
Then the ODE~\eqref{eq:ode_flow} has a unique solution $(X_t)_{t \in [0,T]}$ on $[0, T]$, and for every positive infinitesimal step size $\eta \approx 0$, the hyperfinite Euler recursion~\eqref{eq:euler_scheme} satisfies
\begin{equation}\label{eq:euler_uniform_conv}
\max_{0 \le k\eta \le T} \norm{\ns{X}_{k\eta} - Y_k} \approx 0, \qquad \text{so} \qquad \st\!\left(Y_{\lfloor t/\eta \rfloor}\right) = X_t \quad \text{uniformly on } [0, T].
\end{equation}
\end{theorem}

The conclusion in~\eqref{eq:euler_uniform_conv} compares the whole hyperfinite depth trajectory with the ODE path, rather than only the network state at one chosen layer. With layer $k$ identified with hyperfinite grid time $k\eta \in \mathbb{T}_D$, the approximation error is infinitesimal simultaneously for every internal layer $k$ with $k\eta \le T$, so a single standard ODE path serves as the uniform shadow of the full deep network.

The Euler viewpoint becomes structural in deep residual networks (ResNets)~\cite{he2016deep}.
For comparison, consider an MLP layer of standard width $n \in \N$:
\begin{equation}\label{eq:mlp_layer}
h_{\ell+1} = \frac{1}{\sqrt{n}} W_\ell \varphi(h_\ell), \qquad W_\ell \in \R^{n \times n}.
\end{equation}
In contrast, a (basic) ResNet layer adds a skip connection:
\begin{equation}\label{eq:resnet_layer}
h_{\ell+1} = h_\ell + \frac{1}{\sqrt{n}} W_\ell \varphi(h_\ell).
\end{equation}
The skip term makes it natural to rescale the residual block so that each layer is an infinitesimal update.

A classical deterministic-depth scaling is
\begin{equation}\label{eq:det_resnet}
h_{\ell+1} = h_\ell + \eta \frac{1}{\sqrt{n}} W \ns{\varphi}(h_\ell),
\end{equation}
with a \emph{shared} standard weight matrix $W \in \R^{n \times n}$ across layers (a recurrent-in-depth model). Fix a standard ``layer time'' horizon $\bar{\tau} > 0$ and, for a network of unlimited depth $D \in \hNinf$, set $\eta := \bar{\tau}/D \approx 0$ and $\tau_\ell := \eta\ell \in \mathbb{T}_D = \{0, \eta, 2\eta, \dots, D\eta\}$. Then the update~\eqref{eq:det_resnet} is the exact hyperfinite Euler discretization of
\begin{equation}\label{eq:neural_ode}
\partial_\tau h_\tau = \frac{1}{\sqrt{n}} W \varphi(h_\tau), \qquad \tau \in [0, \bar{\tau}].
\end{equation}
For fixed standard $n$ and $W$, a Lipschitz activation (including ReLU) makes the vector field in~\eqref{eq:neural_ode} globally Lipschitz, so Theorem~\ref{thm:euler_conv} yields $\st(h_{\lfloor \tau/\eta \rfloor}) = h_\tau$ uniformly on $[0, \bar{\tau}]$ at every unlimited depth $D \in \hNinf$. This connects deep ResNets to \emph{Neural ODEs}~\cite{chen2018neural}.

The same viewpoint accommodates non-uniform hyperfinite layer times: varying infinitesimal step sizes $\eta_\ell \approx 0$ and/or layer weights $W_\ell$ can track a time-inhomogeneous ODE whenever the vector field varies S-continuously along the grid.

\subsection{CLT Depth Scaling: From Random Walk to SDE}\label{subsec:clt_resnet}

If each layer contributes random fluctuations of infinitesimal size $\sqrt{\Delta\tau}$ (rather than deterministic size $\Delta\tau$), then a diffusion shadow emerges.

The basic template is Anderson's hyperfinite random walk~\cite{anderson1976nonstandard}, lifting Donsker's random-walk invariance principle~\cite{donsker1951invariance} from~(1.3)--(1.5) in Chapter~1: centered independent increments of size $\sqrt{\Delta\tau}$ over a hyperfinite grid accumulate into an S-continuous trajectory whose standard part is Brownian motion. Allowing state-dependent variance and adding an $\mathcal{O}(\Delta\tau)$ drift gives a Keisler hyperfinite stochastic difference equation~\cite{keisler1984infinitesimal,lindstrom1988invitation} whose shadow solves an It\^o SDE
\begin{equation}\label{eq:general_sde}
dX_\tau = b(X_\tau)\,d\tau + \sigma(X_\tau)\,dB_\tau.
\end{equation}

The same mechanism appears in the \emph{CLT-scaled} ResNet of unlimited depth $D \in \hNinf$ and standard width $n \in \N$, with infinitesimal layer-time step $\Delta\tau := 1/D \approx 0$:
\begin{equation}\label{eq:clt_resnet}
h_{\ell+1} = h_\ell + \sqrt{\Delta\tau}\,\frac{1}{\sqrt{n}} W_\ell \ns{\varphi}(h_\ell),
\end{equation}
where $(W_\ell)_{\ell \ge 0}$ are mutually $\hP$-independent internal $n \times n$ matrices with i.i.d.\ $\hNcal(0, 1)$ entries. Interpret the hyperfinite \emph{layer time} as $\tau_\ell := \ell\Delta\tau$ (so $\ell = \lfloor \tau/\Delta\tau \rfloor$ tracks standard time $\tau$). Since $W_\ell$ is $\hP$-independent of $\mathcal{F}_\ell$, by Transfer the next-layer conditional transition is internally Gaussian:
\begin{equation}\label{eq:clt_resnet_transition}
h_{\ell+1} \mid \mathcal{F}_\ell \sim \hNcal\!\left(h_\ell,\, \Delta\tau\,\frac{1}{n}\norm{\ns{\varphi}(h_\ell)}^2 I_n\right).
\end{equation}
Thus~\eqref{eq:clt_resnet} is, in transition law, an exact Keisler hyperfinite stochastic difference equation corresponding to
\begin{equation}\label{eq:clt_resnet_sde}
dh_\tau = \frac{1}{\sqrt{n}} \norm{\varphi(h_\tau)}\,dB_\tau,
\end{equation}
where $B_\tau$ is a standard $n$-dimensional Brownian motion on the Loeb space $(\Omega, \mathcal{L}(\mathcal{A}), \Loeb)$. For standard $n \in \N$ and a globally Lipschitz activation (including ReLU), the diffusion coefficient $h \mapsto \frac{1}{\sqrt{n}}\norm{\varphi(h)}$ is globally Lipschitz, so Keisler's hyperfinite SDE lifting theorem~\cite{keisler1984infinitesimal} implies that $\tau \mapsto h_{\lfloor \tau/\Delta\tau \rfloor}$ is $\Loeb$-almost surely S-continuous and its shadow $h_\tau = \st(h_{\lfloor \tau/\Delta\tau \rfloor})$ solves~\eqref{eq:clt_resnet_sde}.

When both width $N$ and depth $D$ are allowed to be unlimited in $(\hN)^2$, distinct asymptotic scaling routes correspond to distinct hyperfinite depth--width regimes: the wide galaxy $D/N \approx 0$ (modeling width taken infinite before depth), the deep galaxy $N/D \approx 0$ (depth taken infinite before width), and the proportional galaxy $\st(D/N) = \bar{\tau} \in (0, \infty)$. For the no-bias Gaussian ReLU specialization of~\eqref{eq:clt_resnet}, Hayou and Yang prove that the shadow of the one-neuron preactivation law and empirical covariance flow is independent of the relative magnitudes of $N$ and $D$, provided $\min\{N, D\} \in \hNinf$~\cite[Theorems~1 and~2]{hayou2023width}. In particular, the wide regime $D/N \approx 0$, the deep regime $N/D \approx 0$, and proportional pairs with $\st(D/N) = \bar{\tau}$ all yield the same covariance shadow for these observables. This is a model-specific commutation result: for plain MLPs, proportional and extreme depth--width galaxies can instead disagree. The representative routes in $(\hN)^2$ are summarized in Figure~\ref{fig:depth_width_routes}.

\begin{figure}[htbp]
\centering
\begin{tikzpicture}[
  >=Latex,
  every node/.style={font=\small},
  arrow/.style={->, semithick}
]
  % Left diagram: Gaussian ReLU ResNet
  \begin{scope}[xshift=0cm]
    \node[align=center] at (1.8, 3.95) {Gaussian ReLU ResNet\\{\footnotesize (common covariance shadow)}};

    \node[below left=1pt] at (0, 0) {$(n, d)$};
    \node[above left=1pt] at (0, 2.65) {$(n, D)$};
    \node[below right=1pt] at (3.6, 0) {$(N, d)$};
    \node[above right=1pt] at (3.6, 2.65) {shadow};

    % Bottom arrow (wide galaxy)
    \draw[arrow] (0.05, 0) -- node[below=3pt] {\footnotesize $D/N \approx 0$ (wide)} (3.52, 0);
    % Left arrow (deep galaxy)
    \draw[arrow] (0, 0.05) -- node[left=3pt, align=right] {\footnotesize $N/D \approx 0$\\[-2pt]\footnotesize (deep)} (0, 2.57);
    % Top arrow
    \draw[arrow] (0.08, 2.65) -- node[above=3pt] {\footnotesize $N \in \hNinf$} (3.48, 2.65);
    % Right arrow
    \draw[arrow] (3.6, 0.08) -- node[right=3pt] {\footnotesize $D \in \hNinf$} (3.6, 2.54);
    % Diagonal arrow (proportional galaxy)
    \draw[arrow] (0.06, 0.04) --
      node[pos=0.50, sloped, above=1.5pt] {\footnotesize $\st(D/N) = \bar{\tau}$}
      node[pos=0.50, sloped, below=1.5pt] {\scriptsize (proportional)}
      (3.48, 2.57);
  \end{scope}

  % Right diagram: Plain MLP
  \begin{scope}[xshift=8.2cm]
    \node[align=center] at (1.8, 3.95) {Plain MLP\\{\footnotesize (galaxy dependence can occur)}};

    \node[below left=1pt] at (0, 0) {$(n, d)$};
    \node[above left=1pt] at (0, 2.65) {$(n, D)$};
    \node[below right=1pt] at (3.6, 0) {$(N, d)$};
    \node (neq) at (3.6, 2.65) {$\neq$};

    % Bottom arrow (wide galaxy)
    \draw[arrow] (0.05, 0) -- node[below=3pt] {\footnotesize $D/N \approx 0$ (wide)} (3.52, 0);
    % Left arrow (deep galaxy)
    \draw[arrow] (0, 0.05) -- node[left=3pt, align=right] {\footnotesize $N/D \approx 0$\\[-2pt]\footnotesize (deep)} (0, 2.57);
    % Top arrow (stops short of \neq)
    \draw[arrow] (0.08, 2.65) -- node[above=3pt] {\footnotesize $N \in \hNinf$} (3.25, 2.65);
    % Right arrow (stops short of \neq)
    \draw[arrow] (3.6, 0.08) -- node[right=3pt] {\footnotesize $D \in \hNinf$} (3.6, 2.32);
    % Diagonal arrow (stops short of \neq)
    \draw[arrow] (0.06, 0.04) --
      node[pos=0.48, sloped, above=1.5pt] {\footnotesize $\st(D/N) = \bar{\tau}$}
      node[pos=0.48, sloped, below=1.5pt] {\scriptsize (proportional)}
      (3.28, 2.42);
  \end{scope}
\end{tikzpicture}
\caption{Representative hyperfinite depth--width galaxies in $(\hN)^2$ for initialization observables: the wide galaxy ($D/N \approx 0$), the deep galaxy ($N/D \approx 0$), and the proportional galaxy ($\st(D/N) = \bar{\tau}$). Left: for the no-bias Gaussian ReLU ResNet~\eqref{eq:clt_resnet}, the covariance shadow is independent of the galaxy. Right: for plain MLPs, proportional and sequential galaxies need not commute ($\neq$).}\label{fig:depth_width_routes}
\end{figure}

\section{Proportional Depth--Width Limits at Initialization}\label{sec:prop_linear_init}

To see why joint hyperfinite limits appear for plain deep networks, it is instructive to start with the simplest case: a deep \emph{linear} network at initialization.

Fix standard $n_0 \in \N$ and a nonzero standard input $x \in \R^{n_0}$. Fix unlimited depth and width $D, N \in \hNinf$ with $\st(D/N) = \bar{\tau} \in (0, \infty)$ and infinitesimal depth-time increment $\Delta\tau := 1/N \approx 0$. Let $W_0 \in \hR^{N \times n_0}$ and $(W_\ell)_{\ell \ge 1}$, $W_\ell \in \hR^{N \times N}$, be mutually $\hP$-independent internal matrices with i.i.d.\ $\hNcal(0, 1)$ entries. Set
\begin{equation}\label{eq:deep_linear_net}
h_0 := x, \qquad h_1 := \frac{1}{\sqrt{n_0}} W_0 x, \qquad h_{\ell+1} := \frac{1}{\sqrt{N}} W_\ell h_\ell, \quad \ell \ge 1.
\end{equation}
A hyperfinite depth-$D$ network is the prefix ending at $h_D$. Define the internal normalized squared norms
\begin{equation}\label{eq:scalar_phi_def}
\Phi_0 := \frac{1}{n_0}\norm{x}^2, \qquad \Phi_\ell := \frac{1}{N}\norm{h_\ell}^2, \quad \ell \ge 1.
\end{equation}

\begin{lemma}[Multiplicative Variance Recursion]\label{lem:mult_var_recursion}
For every internal layer $\ell \ge 0$, conditioned on $\mathcal{F}_\ell$, we have by Transfer
\begin{equation}\label{eq:cond_gauss_linear}
h_{\ell+1} \mid \mathcal{F}_\ell \sim \hNcal(0, \Phi_\ell I_N),
\end{equation}
and consequently
\begin{equation}\label{eq:mult_var_step}
\Phi_{\ell+1} = \Phi_\ell \frac{1}{N}\norm{\xi_\ell}^2, \qquad (\xi_\ell)_{\ell \ge 0} \overset{\mathrm{iid}}{\sim} \hNcal(0, I_N),
\end{equation}
where the internal standard Gaussian innovations are $\hP$-independent of the past $\mathcal{F}_\ell$.
\end{lemma}

\begin{proofsketch}
By Transfer of finite-width Gaussian conditioning, for $\ell \ge 1$ write the $i$-th row of $W_\ell$ as $w_{\ell,i}^\top \in \hR^N$. Then
\begin{equation}\label{eq:coord_linear_step}
h_{\ell+1,i} = \frac{1}{\sqrt{N}} w_{\ell,i}^\top h_\ell.
\end{equation}
Conditional on $\mathcal{F}_\ell$, each coordinate is internally centered Gaussian with variance
\begin{equation}\label{eq:coord_linear_var}
\hVar(h_{\ell+1,i} \mid \mathcal{F}_\ell) = \frac{1}{N}\norm{h_\ell}^2 = \Phi_\ell,
\end{equation}
and internal row independence gives the displayed conditional law~\eqref{eq:cond_gauss_linear}. The input step $\ell = 0$ is the same calculation in $\R^{n_0}$ with normalization $n_0^{-1/2}$ and variance $n_0^{-1}\norm{x}^2 = \Phi_0$. Mutual $\hP$-independence of the layer matrices lets us write $h_{\ell+1} = \sqrt{\Phi_\ell}\,\xi_\ell$ with independent internal standard Gaussian innovations $\xi_\ell \sim \hNcal(0, I_N)$. Taking normalized squared norms gives
\begin{equation}\label{eq:phi_xi_norm}
\Phi_{\ell+1} = \frac{1}{N}\norm{h_{\ell+1}}^2 = \Phi_\ell \frac{1}{N}\norm{\xi_\ell}^2,
\end{equation}
which is~\eqref{eq:mult_var_step}.
\end{proofsketch}

Taking internal logarithms turns the hyperfinite product into a hyperfinite sum:
\begin{equation}\label{eq:log_phi_sum}
\ns{\!\log}\!\left(\frac{\Phi_D}{\Phi_0}\right) = \sum_{\ell=0}^{D-1} \ns{\!\log}\!\left(\frac{1}{N}\norm{\xi_\ell}^2\right).
\end{equation}
Since $\frac{1}{N}\norm{\xi_\ell}^2 \approx 1$ with infinitesimal fluctuations of scale $N^{-1/2} = \sqrt{\Delta\tau}$, the sum in~\eqref{eq:log_phi_sum} accumulates to an appreciable random variable precisely when $D$ and $N$ lie in the proportional galaxy ($D/N \in \App$).

\begin{theorem}[Scalar Linear Variance Limit]\label{thm:scalar_linear_var}
For $D, N \in \hNinf$ in the proportional galaxy $\st(D/N) = \bar{\tau} \in (0, \infty)$, the log-variance ratio $\ns{\!\log}(\Phi_D/\Phi_0)$ is $\Loeb$-almost surely limited (so $\Phi_D \in \App$ $\Loeb$-almost surely), and its shadow satisfies
\begin{equation}\label{eq:log_phi_endpoint}
\st\!\left(\ns{\!\log}\!\left(\frac{\Phi_D}{\Phi_0}\right)\right) \sim \mathcal{N}(-\bar{\tau}, 2\bar{\tau}) \qquad \text{under } \Loeb.
\end{equation}
This endpoint law is the equal-hidden-width Gaussian specialization of~\cite[Theorem~1 and Eq.~(9)]{hanin2020products}.

More strongly, define the internal depth-time trajectory $\Phi_\tau^{(N)} := \Phi_{\lfloor N\tau \rfloor}$ with step size $\Delta\tau := 1/N \approx 0$. For every standard $T > 0$, $\tau \mapsto \Phi_\tau^{(N)}$ is $\Loeb$-almost surely limited and S-continuous on $[0, T]$, and its uniform shadow $\Phi_\tau := \st(\Phi_\tau^{(N)})$ is the geometric Brownian motion on $(\Omega, \mathcal{L}(\mathcal{A}), \Loeb)$ solving
\begin{equation}\label{eq:scalar_gbm}
d\Phi_\tau = \sqrt{2}\,\Phi_\tau\,dB_\tau, \qquad \Phi_0 = \frac{1}{n_0}\norm{x}^2.
\end{equation}
The endpoint law in~\eqref{eq:log_phi_endpoint} is a direct consequence of this pathwise shadow.
\end{theorem}

\begin{proofsketch}
By Lemma~\ref{lem:mult_var_recursion}, the internal scalar variance satisfies the multiplicative recursion
\begin{equation}\label{eq:mult_u_def}
\Phi_{\ell+1} = \Phi_\ell U_{\ell,N}, \qquad U_{\ell,N} := \frac{1}{N}\norm{\xi_\ell}^2, \qquad \xi_\ell \sim \hNcal(0, I_N) \text{ i.i.d.}
\end{equation}
Equivalently,
\begin{equation}\label{eq:diff_u_def}
\Phi_{\ell+1} - \Phi_\ell = \Phi_\ell (U_{\ell,N} - 1).
\end{equation}
By Transfer of the $\chi_N^2$ moments, $\norm{\xi_\ell}^2 \sim \ns{\chi}_N^2$ satisfies $\hE[U_{\ell,N} - 1] = 0$, $\hVar(U_{\ell,N} - 1) = 2/N = 2\Delta\tau$, $\hE[(U_{\ell,N} - 1)^4] = \frac{12}{N^2} + \frac{48}{N^3}$, and $\hE[\abs{U_{\ell,N} - 1}^3] = \mathcal{O}(N^{-3/2})$. Conditioning on the forward history $\mathcal{F}_\ell$ with $\Delta\tau = 1/N$ gives
\begin{equation}\label{eq:scalar_cond_moments}
\begin{aligned}
\frac{1}{\Delta\tau}\hE[\Phi_{\ell+1} - \Phi_\ell \mid \mathcal{F}_\ell] &= 0, \\
\frac{1}{\Delta\tau}\hVar(\Phi_{\ell+1} - \Phi_\ell \mid \mathcal{F}_\ell) &= 2\Phi_\ell^2, \\
\frac{1}{\Delta\tau}\hE\!\left[\abs{\Phi_{\ell+1} - \Phi_\ell}^3 \mid \mathcal{F}_\ell\right] &= \mathcal{O}(N^{-1/2})\Phi_\ell^3.
\end{aligned}
\end{equation}
These internal conditional moments depend on the history only through $\Phi_\ell$, matching the zero drift and diffusion covariance $2\Phi^2$ of~\eqref{eq:scalar_gbm} with an infinitesimal third-moment rate on every standard compact set. Defining the Anderson martingale increments $\Delta B_\ell := \frac{1}{\sqrt{2}}(U_{\ell,N} - 1)$ (so $\hE[\Delta B_\ell] = 0$, $\hE[(\Delta B_\ell)^2] = \Delta\tau$, and $B_\tau := \st(\sum_{\ell=0}^{\lfloor N\tau \rfloor - 1} \Delta B_\ell)$ is a standard Brownian motion~\cite{anderson1976nonstandard}), equation~\eqref{eq:diff_u_def} reads $\Phi_{\ell+1} - \Phi_\ell = \sqrt{2}\,\Phi_\ell\,\Delta B_\ell$. Thus Theorem~A.4 gives local S-continuity of $\Phi_\tau^{(N)}$ with shadow solving~\eqref{eq:scalar_gbm} up to bounded-norm stopping. Because both $(\Phi_\ell)_{\ell \ge 0}$ and $(\Phi_\tau)_{\tau \ge 0}$ are nonnegative martingales, Transfer of Doob's maximal inequality gives, for every standard $T > 0$ and $R > 0$,
\begin{equation}\label{eq:doob_maximal}
\hP\!\left(\max_{0 \le \ell \le \lfloor NT \rfloor} \Phi_\ell \ge R\right) \le \frac{\Phi_0}{R}, \qquad \Loeb\!\left(\sup_{0 \le \tau \le T} \Phi_\tau \ge R\right) \le \frac{\Phi_0}{R}.
\end{equation}
Letting standard $R \to \infty$ shows that $\max_{0 \le \ell \le \lfloor NT \rfloor} \Phi_\ell \in \Lim$ $\Loeb$-almost surely, removing the upper stopping barrier and proving global S-continuity on $[0, T]$ with $\st(\Phi_\tau^{(N)}) = \Phi_\tau$.

To see both the endpoint law~\eqref{eq:log_phi_endpoint} and the It\^o drift directly at the hyperfinite level, observe that for every standard $\varepsilon > 0$, Markov's inequality yields $\hP(\max_{0 \le \ell < D} \abs{U_{\ell,N} - 1} > \varepsilon) \le D\varepsilon^{-4}\hE[(U_{0,N} - 1)^4] = \mathcal{O}(D/N^2) \approx 0$, so $\max_{0 \le \ell < D} \abs{U_{\ell,N} - 1} \approx 0$ $\Loeb$-almost surely. Summing the nonstandard second-order Taylor expansion $\ns{\!\log}(1 + (U_{\ell,N} - 1)) = (U_{\ell,N} - 1) - \frac{1}{2}(U_{\ell,N} - 1)^2 + \mathcal{O}(\abs{U_{\ell,N} - 1}^3)$ over $\ell = 0, \dots, D-1$ in~\eqref{eq:log_phi_sum} gives
\begin{equation}\label{eq:ito_log_phi}
\begin{aligned}
\ns{\!\log}\!\left(\frac{\Phi_D}{\Phi_0}\right) &= \sum_{\ell=0}^{D-1} (U_{\ell,N} - 1) - \frac{1}{2}\sum_{\ell=0}^{D-1} (U_{\ell,N} - 1)^2 + \sum_{\ell=0}^{D-1} \mathcal{O}\!\left(\abs{U_{\ell,N} - 1}^3\right) \\
&= \sqrt{2}\sum_{\ell=0}^{D-1} \Delta B_\ell - \frac{D}{N} + \Inf \approx \sqrt{2}\,B_{\bar{\tau}} - \bar{\tau} \qquad (\Loeb\text{-a.s.}),
\end{aligned}
\end{equation}
because $-\frac{1}{2}\sum_{\ell=0}^{D-1} (U_{\ell,N} - 1)^2$ has internal mean $-\sum_{\ell=0}^{D-1} \frac{1}{2}\cdot\frac{2}{N} = -\frac{D}{N} \approx -\bar{\tau}$ and internal variance $\mathcal{O}(D/N^2) \approx 0$, while $\sum_{\ell=0}^{D-1} \hE[\abs{U_{\ell,N} - 1}^3] = \mathcal{O}(D N^{-3/2}) \approx 0$. Equivalently, at the shadow level, It\^o's lemma $d\log\Phi_\tau = \frac{1}{\Phi_\tau}\,d\Phi_\tau - \frac{1}{2\Phi_\tau^2}(d\Phi_\tau)^2 = -d\tau + \sqrt{2}\,dB_\tau$ yields $\log(\Phi_{\bar{\tau}}/\Phi_0) = -\bar{\tau} + \sqrt{2}\,B_{\bar{\tau}} \sim \mathcal{N}(-\bar{\tau}, 2\bar{\tau})$, proving~\eqref{eq:log_phi_endpoint}.
\end{proofsketch}

The term \emph{geometric} in~\eqref{eq:scalar_gbm} refers to its multiplicative noise: randomness scales with $\Phi_\tau$ itself. This is the probabilistic signature of hyperfinite matrix products, in which each layer randomly stretches or shrinks the current state.

Even in this simplest linear setting, the proportional galaxy already produces \emph{non-Gaussian} marginal shadows. For any fixed coordinate $i$, $h_{D,i} \mid \mathcal{F}_{D-1} \sim \hNcal(0, \Phi_{D-1})$. Since S-continuity gives $\st(\Phi_{D-1}) = \Phi_{\bar{\tau}}$ when $\st(D/N) = \bar{\tau}$, the standard-part marginal $\st(h_{D,i})$ under $\Loeb$ is a Gaussian scale mixture with random mixing variance $\Phi_{\bar{\tau}}$. This is a key motivation for working in the proportional depth--width galaxy.

Here ``stable'' means appreciable in a two-sided probabilistic sense. For a positive internal random variable $X > 0$, the condition $\Loeb(X \in \Lim) = 1$ (the nonstandard counterpart of $X_k = \mathcal{O}_\Prob(1)$) only rules out escape to unlimited values in $\hR \setminus \Lim$ and still permits infinitesimal collapse $X \approx 0$. Two-sided probabilistic stability ($X_k = \Theta_\Prob(1)$) translates cleanly to $X$ lying in the appreciable group $\App = \Lim \setminus \Inf$ with $\Loeb$-probability $1$ (equivalently, both $X$ and $X^{-1}$ lie in $\Lim$ $\Loeb$-almost surely), so typical realizations have positive finite shadows $\st(X) \in (0, \infty)$.

The three hyperfinite depth--width galaxies behave differently with respect to $\App$. In the wide galaxy ($D/N \approx 0$, such as $N \in \hNinf$ at standard depth $d \in \N$), layerwise fluctuations average away in~\eqref{eq:ito_log_phi} and $\Phi_D \approx \Phi_0 \in \App$ $\Loeb$-almost surely, collapsing to the deterministic appreciable value $\Phi_0$. In the deep galaxy ($N/D \approx 0$, such as $D \in \hNinf$ at standard width $n \in \N$), repeated multiplication instead drives $\ns{\!\log}(\Phi_D/\Phi_0)$ to $-\infty$, so $\Phi_D \approx 0$ $\Loeb$-almost surely even though the internal expectation $\hE[\Phi_D] = \Phi_0$ remains constant for every $D$. Thus $\Loeb(\Phi_D \in \Lim) = 1$ whereas $\Loeb(\Phi_D \in \App) = 0$: an internal event of infinitesimal probability carrying hyper-large realizations preserves $\hE[\Phi_D]$ (a failure of S-integrability), while $\Loeb$-almost every realization collapses into $\Inf$. Under proportional scaling $\st(D/N) = \bar{\tau} \in (0, \infty)$, by contrast, $\st(\Phi_D) = \Phi_{\bar{\tau}} \in (0, \infty)$ $\Loeb$-almost surely, so $\Phi_D$ remains in $\App$ with $\Loeb$-probability $1$ while retaining genuine randomness~\cite[\S1.1, Theorem~1, and Eq.~(9)]{hanin2020products}.

Although the discussion so far concerns initialization, proportional depth--width scaling also changes the NTK picture for fully connected ReLU networks. For equal hidden width $N$, Hanin and Nica show that the on-diagonal NTK at initialization need not concentrate to a deterministic shadow when $\st(D/N) > 0$~\cite[Theorem~1]{hanin2020finite}. Their training-time result is narrower: for square loss and a single-example SGD step, they compute a potentially appreciable mean first update of the same kernel entry. This motivates, but does not prove, their conjectured ``weak feature learning'' regime~\cite[Theorem~2 and \S3.2]{hanin2020finite}.

\section{Multiple Data Points: A Covariance SDE}\label{sec:multi_cov_sde}

We write $\Sym(m)$ for the vector space of real symmetric $m \times m$ matrices (with internal counterpart $\ns{\Sym}(m)$ over $\hR$). The internal empirical covariance matrices below are positive semidefinite; they need not be strictly positive definite.

Now fix standard $m \ge 2$ and propagate standard inputs $x^1, \dots, x^m \in \R^{n_0}$ through the same hyperfinite random linear network~\eqref{eq:deep_linear_net} of width $N \in \hNinf$. Define the input and internal hidden-layer Gram matrices by
\begin{equation}\label{eq:multi_linear_gram}
\begin{aligned}
\Phi_0^{\alpha\beta} &:= \frac{1}{n_0}\ip{x^\alpha}{x^\beta}, \qquad \Phi_\ell^{\alpha\beta} := \frac{1}{N}\ip{h_\ell^\alpha}{h_\ell^\beta}, \quad \ell \ge 1, \\
\Phi_\ell &:= \left[\Phi_\ell^{\alpha\beta}\right]_{\alpha,\beta=1}^m \in \ns{\Sym}(m), \quad \ell \ge 0.
\end{aligned}
\end{equation}
Let $\mathcal{F}_0 := \{\varnothing, \Omega\}$ be the trivial algebra, and for $\ell \ge 1$ define the internal algebra
\begin{equation}\label{eq:internal_filtration}
\mathcal{F}_\ell := \ns{\sigma}\!\left(\left\{h_k^\alpha\right\}_{1 \le k \le \ell,\; \alpha \in [m]}\right).
\end{equation}
Thus $\mathcal{F}_\ell$ records the forward hidden-state history through layer $\ell$. Conditioned on $\mathcal{F}_\ell$, by Transfer $\{h_{\ell+1}^\alpha\}_{\alpha=1}^m$ is internally jointly Gaussian with covariance $\Phi_\ell \otimes I_N$. As a result, each entry of the centered increment $\Phi_{\ell+1} - \Phi_\ell$ is a hyperfinite empirical average of conditionally i.i.d.\ centered terms. For $1 \le \alpha \le \beta \le m$, write
\begin{equation}\label{eq:gram_inc_decomp}
\Phi_{\ell+1}^{\alpha\beta} = \Phi_\ell^{\alpha\beta} + \frac{1}{\sqrt{N}} R_\ell^{\alpha\beta}, \qquad \hE\!\left[R_\ell^{\alpha\beta} \mid \mathcal{F}_\ell\right] = 0,
\end{equation}
where
\begin{equation}\label{eq:gram_fluct_R}
R_\ell^{\alpha\beta} := \frac{1}{\sqrt{N}} \sum_{i=1}^N \left(h_{\ell+1,i}^\alpha h_{\ell+1,i}^\beta - \Phi_\ell^{\alpha\beta}\right).
\end{equation}

The next step is to compute the internal conditional covariance between these fluctuation entries.

\begin{proposition}[Isserlis Covariance for the Gram Increment]\label{prop:isserlis_gram}
Let indices satisfy $\alpha \le \beta$ and $\gamma \le \delta$ (to avoid duplication from symmetry), and let $R_\ell^{\alpha\beta}$ be the centered internal Gram fluctuation defined in~\eqref{eq:gram_fluct_R}. Then, conditional on $\mathcal{F}_\ell$, the neuron-wise internal vectors
\begin{equation}\label{eq:neuron_vec_linear}
h_{\ell+1,i} := \left(h_{\ell+1,i}^1, \dots, h_{\ell+1,i}^m\right) \in \hR^m, \qquad i = 1, \dots, N,
\end{equation}
are internally i.i.d.\ Gaussian with law $\hNcal(0, \Phi_\ell)$, and by Transfer we have the \emph{exact} internal identities
\begin{equation}\label{eq:isserlis_cov_identities}
\hE\!\left[R_\ell^{\alpha\beta} \mid \mathcal{F}_\ell\right] = 0, \qquad \hCov\!\left(R_\ell^{\alpha\beta}, R_\ell^{\gamma\delta} \mid \mathcal{F}_\ell\right) = \Phi_\ell^{\alpha\gamma}\Phi_\ell^{\beta\delta} + \Phi_\ell^{\alpha\delta}\Phi_\ell^{\beta\gamma}.
\end{equation}
Equivalently, in the flattened upper-triangular coordinates on $\Sym(m)$, the conditional covariance matrix is $\Sigma(\Phi_\ell)$ with entries
\begin{equation}\label{eq:sigma_tensor_def}
\Sigma_{\alpha\beta,\gamma\delta}(\Phi) := \Phi^{\alpha\gamma}\Phi^{\beta\delta} + \Phi^{\alpha\delta}\Phi^{\beta\gamma}.
\end{equation}
\end{proposition}

\begin{proof}
Fix $\ell$ and condition on $\mathcal{F}_\ell$. For each $i \in \{1, \dots, N\}$, set
\begin{equation}\label{eq:single_neuron_R}
R_{\ell,i}^{\alpha\beta} := h_{\ell+1,i}^\alpha h_{\ell+1,i}^\beta - \Phi_\ell^{\alpha\beta}, \qquad \text{so that} \qquad R_\ell^{\alpha\beta} = \frac{1}{\sqrt{N}} \sum_{i=1}^N R_{\ell,i}^{\alpha\beta}.
\end{equation}
Since $(h_{\ell+1,i})_{i=1}^N$ are conditionally internally i.i.d.\ $\hNcal(0, \Phi_\ell)$, the centered variables $(R_{\ell,i}^{\alpha\beta})_{i=1}^N$ are conditionally i.i.d.\ as well, and therefore by Transfer
\begin{equation}\label{eq:cov_reduce_i1}
\hCov\!\left(R_\ell^{\alpha\beta}, R_\ell^{\gamma\delta} \mid \mathcal{F}_\ell\right) = \hE\!\left[R_{\ell,1}^{\alpha\beta} R_{\ell,1}^{\gamma\delta} \mid \mathcal{F}_\ell\right].
\end{equation}
Expanding and using $\hE[h_{\ell+1,1}^\alpha h_{\ell+1,1}^\beta \mid \mathcal{F}_\ell] = \Phi_\ell^{\alpha\beta}$ gives
\begin{equation}\label{eq:fourth_moment_centered}
\hE\!\left[R_{\ell,1}^{\alpha\beta} R_{\ell,1}^{\gamma\delta} \mid \mathcal{F}_\ell\right] = \hE\!\left[h_{\ell+1,1}^\alpha h_{\ell+1,1}^\beta h_{\ell+1,1}^\gamma h_{\ell+1,1}^\delta \mid \mathcal{F}_\ell\right] - \Phi_\ell^{\alpha\beta}\Phi_\ell^{\gamma\delta}.
\end{equation}
By Transfer of Isserlis' theorem (Wick's formula) to the centered internal Gaussian vector $h_{\ell+1,1}$,
\begin{equation}\label{eq:wick_four_point}
\hE\!\left[h_{\ell+1,1}^\alpha h_{\ell+1,1}^\beta h_{\ell+1,1}^\gamma h_{\ell+1,1}^\delta \mid \mathcal{F}_\ell\right] = \Phi_\ell^{\alpha\beta}\Phi_\ell^{\gamma\delta} + \Phi_\ell^{\alpha\gamma}\Phi_\ell^{\beta\delta} + \Phi_\ell^{\alpha\delta}\Phi_\ell^{\beta\gamma},
\end{equation}
and subtracting the first term yields the claimed covariance~\eqref{eq:isserlis_cov_identities}.
\end{proof}

By Proposition~\ref{prop:isserlis_gram}, with $\Delta\tau := 1/N \approx 0$ and $\bar{m} := m(m+1)/2$, we may factor the flattened fluctuation as $R_\ell = \Sigma(\Phi_\ell)^{1/2}\xi_\ell$ with $\hE[\xi_\ell \mid \mathcal{F}_\ell] = 0$ and $\hCov(\xi_\ell \mid \mathcal{F}_\ell) = I_{\bar{m}}$ (adjoining independent $\hNcal(0,1)$ coordinates along any null directions of $\Sigma(\Phi_\ell)$), so one Gram increment takes the exact Keisler martingale form $\Delta\Phi_\ell := \Phi_{\ell+1} - \Phi_\ell = \Sigma(\Phi_\ell)^{1/2}\Delta B_\ell$ with $\Delta B_\ell := N^{-1/2}\xi_\ell$. Thus, on the hyperfinite layer-time scale $\tau_\ell = \ell/N$, fluctuations across $\mathcal{O}(N)$ layers accumulate at appreciable order. The following theorem formalizes this pathwise shadow.

\begin{theorem}[Linear Covariance SDE Limit {\cite[Theorem~3.2 and Appendix~C.1]{li2022neural}}]\label{thm:linear_cov_sde}
Fix $N \in \hNinf$ and define the internal piecewise-constant process $\Phi_\tau^{(N)} := \Phi_{\lfloor N\tau \rfloor}$ for $\tau \ge 0$. Let $\bar{m} := m(m+1)/2$. After flattening $\Phi$ into its $\bar{m}$ distinct upper-triangular entries, for every standard $\bar{\tau} > 0$, the trajectory $\tau \mapsto \Phi_\tau^{(N)}$ is $\Loeb$-almost surely limited and S-continuous on $[0, \bar{\tau}]$, and its uniform shadow $\Phi_\tau := \st(\Phi_\tau^{(N)})$ solves the covariance SDE on $(\Omega, \mathcal{L}(\mathcal{A}), \Loeb)$:
\begin{equation}\label{eq:linear_cov_sde}
d\Phi_\tau = \Sigma(\Phi_\tau)^{1/2}\,dB_\tau, \qquad \Phi_0 = \left[\frac{1}{n_0}\ip{x^\alpha}{x^\beta}\right]_{\alpha,\beta=1}^m.
\end{equation}
Here $B_\tau = \st(\sum_{\ell=0}^{\lfloor N\tau \rfloor - 1} \Delta B_\ell)$ is a standard Brownian motion in $\R^{\bar{m}}$, and $\Sigma$ is the covariance matrix in~\eqref{eq:sigma_tensor_def}.
\end{theorem}

\begin{proofsketch}
By~\eqref{eq:gram_inc_decomp} and Proposition~\ref{prop:isserlis_gram}, with $\Delta\tau = 1/N$, one layer has internal conditional mean zero and conditional covariance $\Delta\tau\,\Sigma(\Phi_\ell)$, so $\Delta\Phi_\ell = \Sigma(\Phi_\ell)^{1/2}\Delta B_\ell$ with $\Delta B_\ell = N^{-1/2}\xi_\ell$. Transferred Gaussian moment bounds give $\frac{1}{\Delta\tau}\hE[\norm{\Delta\Phi_\ell}^3 \mid \mathcal{F}_\ell] = \mathcal{O}(N^{-1/2})\norm{\Phi_\ell}_F^3 \approx 0$ on limited states, verifying the third-moment hypothesis of Theorem~A.4. Because these conditional moments depend on the history only through $\Phi_\ell$, they furnish the Markov-state coefficients required by Theorem~A.4. Finally, applying Doob's maximal inequality~\eqref{eq:doob_maximal} to each nonnegative diagonal martingale $\Phi_\ell^{\alpha\alpha}$ and Cauchy--Schwarz $\abs{\Phi_\ell^{\alpha\beta}}^2 \le \Phi_\ell^{\alpha\alpha}\Phi_\ell^{\beta\beta}$ to the off-diagonal entries ensures $\max_{0 \le \ell \le \lfloor N\bar{\tau} \rfloor} \norm{\Phi_\ell}_F \in \Lim$ $\Loeb$-almost surely, yielding global S-continuity and the shadow SDE~\eqref{eq:linear_cov_sde} via Theorem~A.4.
\end{proofsketch}

When $m = 1$,~\eqref{eq:sigma_tensor_def} gives $\Sigma(\Phi) = 2\Phi^2$, so Theorem~\ref{thm:linear_cov_sde} reduces to the geometric Brownian motion in Theorem~\ref{thm:scalar_linear_var}. If the standard inputs are linearly dependent, those linear relations are preserved identically by every internal linear layer. One may therefore apply the theorem to a maximal linearly independent subset and reconstruct the full Gram shadow from it.

The same covariance tensor has a geometric interpretation: on $\SPD(m)$, $\Sigma(\Phi)$ is the cometric, or inverse metric tensor, associated with the affine-invariant Riemannian metric. See Section~6.2.2 for this metric interpretation and Section~6.2.3 for the resulting eigenvalue dynamics.

\section{Nonlinearities, Correlation Collapse, and Shaping}\label{sec:shaping}

For nonlinear activations, the hyperfinite depth shadow depends strongly on normalization and on how quickly correlations are driven toward a fixed point.

Consider first the ReLU activation $\varphi(u) = \max\{u, 0\}$. If $g \sim \mathcal{N}(0, \sigma^2)$ then $\E[\varphi(g)^2] = \sigma^2/2$. Thus, for the naive hyperfinite ReLU MLP
\begin{equation}\label{eq:naive_relu_mlp}
h_{\ell+1} = \frac{1}{\sqrt{N}} W_\ell \ns{\varphi}(h_\ell),
\end{equation}
the variance contracts by a factor $1/2$ per layer in conditional expectation. The standard fix is Kaiming normalization~\cite{he2015delving}: choose $c > 0$ so that $c^{-1} = \E[\varphi(\omega)^2]$ for $\omega \sim \mathcal{N}(0, 1)$, and scale the layer by $\sqrt{c}$. For ReLU, $c = 2$.

\begin{proposition}[ReLU SDE Limit for One Sample]\label{prop:relu_sde_one_sample}
Fix standard $n_0 \in \N$ and $h_0 := x \in \R^{n_0} \setminus \{0\}$. For $D, N \in \hNinf$ with $\st(D/N) = \bar{\tau} \in (0, \infty)$, consider the Kaiming-normalized hyperfinite ReLU network
\begin{equation}\label{eq:kaiming_relu_net}
h_1 = \frac{1}{\sqrt{n_0}} W_0 x, \qquad h_{\ell+1} = \sqrt{\frac{c}{N}} W_\ell \ns{\varphi}(h_\ell), \qquad \ell \ge 1, \qquad c = 2.
\end{equation}
Here $W_0 \in \hR^{N \times n_0}$ and $W_\ell \in \hR^{N \times N}$ for $\ell \ge 1$ are mutually $\hP$-independent internal matrices with i.i.d.\ $\hNcal(0, 1)$ entries. Define the (normalized) internal squared post-activation norm
\begin{equation}\label{eq:kaiming_relu_phi}
\Phi_0 := \frac{1}{n_0}\norm{x}_2^2, \qquad \Phi_\ell := \frac{c}{N}\norm{\ns{\varphi}(h_\ell)}_2^2, \qquad \ell \ge 1.
\end{equation}
For $\tau \in [0, \bar{\tau}]$, set $\Phi_\tau^{(N)} := \Phi_{\lfloor \tau N \rfloor}$. Then $\tau \mapsto \Phi_\tau^{(N)}$ is $\Loeb$-almost surely limited and S-continuous on $[0, \bar{\tau}]$ (see Appendix~A.3), and its uniform shadow $\Phi_\tau := \st(\Phi_\tau^{(N)})$ is the geometric Brownian motion on $(\Omega, \mathcal{L}(\mathcal{A}), \Loeb)$ solving
\begin{equation}\label{eq:relu_gbm}
d\Phi_\tau = \sqrt{5}\,\Phi_\tau\,dB_\tau,
\end{equation}
with the same initial condition $\Phi_0 = \frac{1}{n_0}\norm{x}_2^2$. At the terminal hyperfinite depth $D$, $\Phi_D \in \App$ $\Loeb$-almost surely and
\begin{equation}\label{eq:relu_terminal_log}
\st\!\left(\ns{\!\log}\!\left(\frac{\Phi_D}{\Phi_0}\right)\right) \sim \mathcal{N}\!\left(-\frac{5}{2}\bar{\tau},\, 5\bar{\tau}\right) \qquad \text{under } \Loeb.
\end{equation}
\end{proposition}

The proof is deferred to Exercise~\ref{ex:relu_sde_one_sample}.

For one input, Kaiming normalization gives $\hE[\Phi_{\ell+1} \mid \mathcal{F}_\ell] = \Phi_\ell$; the remaining scale-$N^{-1/2}$ hyperfinite fluctuation accumulates into the SDE of Proposition~\ref{prop:relu_sde_one_sample}. This controls only the diagonal of the multi-input Gram matrix; the off-diagonal correlation update remains the obstruction. For two inputs with nonzero postactivation norms, define the internal scaled postactivation Gram entry and its correlation at layer $\ell$ by
\begin{equation}\label{eq:relu_gram_corr}
\Phi_\ell^{\alpha\beta} := \frac{c}{N}\ip{\ns{\varphi}(h_\ell^\alpha)}{\ns{\varphi}(h_\ell^\beta)}, \qquad \rho_\ell^{\alpha\beta} := \frac{\Phi_\ell^{\alpha\beta}}{\sqrt{\Phi_\ell^{\alpha\alpha}\Phi_\ell^{\beta\beta}}} \in \ns{[-1, 1]}.
\end{equation}
In the wide galaxy ($N \in \hNinf$ at standard depth $\ell \in \N$), the Cho--Saul formula in Proposition~2.5, applied with unit variances, gives the shadow ReLU correlation recursion $\st(\rho_{\ell+1}^{\alpha\beta}) = \mathcal{T}(\st(\rho_\ell^{\alpha\beta}))$:
\begin{equation}\label{eq:cho_saul_corr_map}
\rho_{\ell+1}^{\alpha\beta} \approx \mathcal{T}\!\left(\rho_\ell^{\alpha\beta}\right), \qquad \mathcal{T}(\rho) := \frac{\E[\varphi(w_1)\varphi(w_2)]}{\E[\varphi(g)^2]} = \frac{1}{\pi}\left(\sqrt{1 - \rho^2} + (\pi - \arccos\rho)\,\rho\right),
\end{equation}
where $g \sim \mathcal{N}(0, 1)$ and $(w_1, w_2)$ is a jointly standard Gaussian pair with correlation $\rho$. For $-1 \le \rho < 1$, one has $\mathcal{T}(\rho) > \rho$; hence $1$ is the unique fixed point in $[-1, 1]$, and every standard correlation recursion starting below $1$ converges to $1$ (so over an unlimited number of layers, correlations collapse into the monad of $1$).

\begin{remark}[Representation and Kernel Collapse]\label{rem:kernel_collapse}
As an off-diagonal correlation enters the monad of $1$ ($\rho_\ell^{\alpha\beta} \approx 1$), the corresponding normalized postactivation representations align up to infinitesimals and their normalized kernel shadow collapses to its diagonal value. This is a forward representation/kernel-collapse statement; it does not by itself imply that backpropagated gradients vanish.
\end{remark}

Away from $1$, the map $\mathcal{T}$ in~\eqref{eq:cho_saul_corr_map} differs from the identity by an appreciable ($\App$-scale) amount, so the unshaped off-diagonal update is not on the infinitesimal drift scale $\Delta\tau = 1/N$ needed over $D \asymp N$ layers. Activation shaping repairs this scale mismatch: the full normalized kernel map becomes the identity plus an $\mathcal{O}(\Delta\tau)$ bias while the empirical Gram fluctuation remains of scale $\sqrt{\Delta\tau} = N^{-1/2}$. This viewpoint is closely related to \emph{deep kernel shaping} and \emph{tailored rectifiers}~\cite{martens2021rapid,zhang2022deep}, and it underlies the covariance SDE limits of~\cite[Sections~2.3--3.2]{li2022neural}.

For concreteness, fix standard inputs $x^1, \dots, x^m \in \R^{n_0}$ and $D, N \in \hNinf$ with $\st(D/N) = \bar{\tau} \in (0, \infty)$, and consider the Kaiming-normalized shaped hyperfinite network
\begin{equation}\label{eq:shaped_net}
h_1^\alpha = \frac{1}{\sqrt{n_0}} W_0 x^\alpha, \qquad h_{\ell+1}^\alpha = \sqrt{\frac{c}{N}} W_\ell \varphi_s(h_\ell^\alpha), \qquad 1 \le \ell \le D - 1,
\end{equation}
where $W_0 \in \hR^{N \times n_0}$ and $W_\ell \in \hR^{N \times N}$ are mutually $\hP$-independent with i.i.d.\ $\hNcal(0, 1)$ entries, and $c = c(N) \in \hR_{>0}$ is chosen so that $\hE[\varphi_s(g)^2] = c^{-1}$ for $g \sim \hNcal(0, 1)$. Define the input boundary and the internal scaled postactivation Gram matrices by
\begin{equation}\label{eq:shaped_gram}
\begin{aligned}
\Phi_0^{\alpha\beta} &:= \frac{1}{n_0}\ip{x^\alpha}{x^\beta}, \\
\Phi_\ell^{\alpha\beta} &:= \frac{c}{N}\ip{\varphi_s(h_\ell^\alpha)}{\varphi_s(h_\ell^\beta)}, \quad 1 \le \ell \le D, \qquad \Phi_\ell := \left(\Phi_\ell^{\alpha\beta}\right)_{\alpha,\beta=1}^m \in \ns{\Sym}(m).
\end{aligned}
\end{equation}
Thus $\Phi_\ell$ for $\ell \ge 1$ is the internal postactivation Gram matrix at layer $\ell$. Every empirical $\Phi_\ell$ is internally positive semidefinite. We assume that the standard data Gram matrix $\Phi_0$ is strictly positive definite, although the one-step calculations below remain valid for positive-semidefinite covariances. The dependence of the hyperfinite chain $\Phi_\ell$ on $N$ is suppressed. With $\mathcal{F}_\ell$ as in~\eqref{eq:internal_filtration}, conditional on $\mathcal{F}_\ell$ the vectors $h_{\ell+1,i} = (h_{\ell+1,i}^1, \dots, h_{\ell+1,i}^m)$, $i = 1, \dots, N$, are internally i.i.d.\ $\hNcal(0, \Phi_\ell)$ for $0 \le \ell \le D - 1$.

\paragraph{ReLU-like shaping.}
For standard $p > 0$ and $c_+, c_- \in \R$, define the internal piecewise-linear activation
\begin{equation}\label{eq:relu_like_shaping}
\varphi_s(x) = s_+ \max\{x, 0\} + s_- \min\{x, 0\}, \qquad s_\pm = 1 + \frac{c_\pm}{N^p}.
\end{equation}
At the critical exponent $p = \frac{1}{2}$ (so $s_\pm = 1 + c_\pm / \sqrt{N} \approx 1$), this is~\cite[Definition~3.1]{li2022neural}; the general exponent family is analyzed in~\cite[Proposition~3.4]{li2022neural}. The exponent $p = \frac{1}{2}$ is critical because the deterministic bias of the normalized kernel is then $\mathcal{O}(\Delta\tau)$ (with $\Delta\tau := 1/N$) and accumulates to an appreciable drift over $D \asymp N$ layers.

\paragraph{Smooth shaping.}
Assume standard $\varphi \in C^4(\R)$ satisfies $\varphi(0) = 0$, $\varphi'(0) = 1$, and $\abs{\varphi^{(4)}(x)} \le C(1 + \abs{x}^q)$ for some standard $C, q > 0$, as in~\cite[Assumption~3.5]{li2022neural}. For a standard constant $a > 0$, use the unlimited rescaling
\begin{equation}\label{eq:smooth_shaping}
\varphi_s(x) := s\,\ns{\varphi}(x/s), \qquad s = a\sqrt{N}, \qquad c^{-1} = \hE[\varphi_s(g)^2], \quad g \sim \hNcal(0, 1).
\end{equation}
Then for every limited $x \in \Lim$, $x/s \approx 0$ so $\varphi_s(x) \approx x$, with the transferred fourth-derivative bound controlling the nonstandard Taylor remainder in internal expectation, making the kernel map an infinitesimal perturbation of the identity~\cite[Definition~3.6]{li2022neural}.

\paragraph{Conditional remainder notation.}
Let $\Delta\tau := 1/N \approx 0$. For an internal matrix- or vector-valued random variable $X$, possibly depending on $\ell$ and $N$, we write $X = \tilde{o}_\ell(\Delta\tau)$ whenever (at every limited state $\Phi_\ell$, uniformly on standard compact subsets)
\begin{equation}\label{eq:cond_o_tilde}
\frac{\norm{\hE[X \mid \mathcal{F}_\ell]}_F}{\Delta\tau} \approx 0, \qquad \frac{\hE\!\left[\norm{X - \hE[X \mid \mathcal{F}_\ell]}_F^2 \;\middle|\; \mathcal{F}_\ell\right]}{\Delta\tau} \approx 0.
\end{equation}
The subscript in $\tilde{o}_\ell(\Delta\tau)$ records conditioning on $\mathcal{F}_\ell$. Thus its internal conditional mean is infinitesimal relative to $\Delta\tau = 1/N$, while its centered part is infinitesimal relative to $\sqrt{\Delta\tau} = N^{-1/2}$ in conditional internal $L^2$; this does not assert $X/\Delta\tau \approxP 0$. The conditioning rules out predictable one-step errors that disappear only after averaging over the current state yet could accumulate over $D \asymp N$ layers.

\begin{theorem}[Keisler--Euler--Maruyama Increment Expansion under Critical Shaping]\label{thm:kem_increment}
Assume either ReLU-like shaping~\eqref{eq:relu_like_shaping} at the critical exponent $p = \frac{1}{2}$ ($s_\pm = 1 + c_\pm/\sqrt{N}$), or smooth shaping~\eqref{eq:smooth_shaping} with $s = a\sqrt{N}$, and set $\Delta\tau := 1/N \approx 0$. After flattening the upper-triangular entries as coordinates on $\Sym(m)$, whenever $\Phi_\ell$ is limited one can write
\begin{equation}\label{eq:kem_step}
\Phi_{\ell+1} - \Phi_\ell = b(\Phi_\ell)\,\Delta\tau + \Sigma(\Phi_\ell)^{1/2}\Delta B_\ell + \tilde{o}_\ell(\Delta\tau),
\end{equation}
where $\tilde{o}_\ell(\Delta\tau)$ is understood in the sense of~\eqref{eq:cond_o_tilde} and $\Delta B_\ell := \sqrt{\Delta\tau}\,\xi_\ell = N^{-1/2}\xi_\ell$, with internal noise vector $\xi_\ell$ satisfying
\begin{equation}\label{eq:xi_cond_cov}
\hE[\xi_\ell \mid \mathcal{F}_\ell] = 0, \qquad \hCov(\xi_\ell \mid \mathcal{F}_\ell) = I_{\bar{m}}, \qquad \bar{m} := \frac{m(m+1)}{2},
\end{equation}
and $\Sigma$ the covariance matrix in~\eqref{eq:sigma_tensor_def}. The drift $b$ (extended naturally to internal $\Phi_\ell$) is given as follows.
\begin{itemize}
\item In the ReLU-like shaping scheme, one has the entrywise formula
\begin{equation}\label{eq:drift_relu_like}
\begin{aligned}
b^{\alpha\beta}(\Phi) &= \sqrt{\Phi^{\alpha\alpha}\Phi^{\beta\beta}}\,\nu(\rho^{\alpha\beta}), \qquad \rho^{\alpha\beta} := \frac{\Phi^{\alpha\beta}}{\sqrt{\Phi^{\alpha\alpha}\Phi^{\beta\beta}}}, \\
\nu(\rho) &:= \frac{(c_+ - c_-)^2}{2\pi}\left(\sqrt{1 - \rho^2} - \rho\arccos(\rho)\right).
\end{aligned}
\end{equation}
\item In the smooth-shaping scheme, the drift admits the entrywise form
\begin{equation}\label{eq:drift_smooth}
b^{\alpha\beta}(\Phi) = \frac{\varphi''(0)^2}{4a^2}\Bigl(\Phi^{\alpha\alpha}\Phi^{\beta\beta} + \Phi^{\alpha\beta}(2\Phi^{\alpha\beta} - 3)\Bigr) + \frac{\varphi'''(0)}{2a^2}\Phi^{\alpha\beta}\Bigl(\Phi^{\alpha\alpha} + \Phi^{\beta\beta} - 2\Bigr).
\end{equation}
\end{itemize}
The ReLU-like formula is understood by continuity if a diagonal covariance entry vanishes.
\end{theorem}

\begin{proof}
Fix $\ell$ with limited $\Phi_\ell$, and condition on $\mathcal{F}_\ell$. Since $W_\ell$ has internally i.i.d.\ $\hNcal(0,1)$ rows, by Transfer the neuron-wise preactivation vectors
\begin{equation}\label{eq:shaped_neuron_vec}
h_{\ell+1,i} := \left(h_{\ell+1,i}^1, \dots, h_{\ell+1,i}^m\right) \in \hR^m, \qquad i = 1, \dots, N,
\end{equation}
are conditionally internally i.i.d.\ $\hNcal(0, \Phi_\ell)$.

\medskip\noindent\emph{Drift calculation.}
Let $w \mid \mathcal{F}_\ell \sim \hNcal(0, \Phi_\ell)$. In the ReLU-like case, set $\rho_\ell^{\alpha\beta} := \Phi_\ell^{\alpha\beta}/\sqrt{\Phi_\ell^{\alpha\alpha}\Phi_\ell^{\beta\beta}}$ and let $(z_1, z_2)$ be an internal standard Gaussian pair with correlation $\rho_\ell^{\alpha\beta}$. By Transfer of positive homogeneity and the arc-cosine calculation~\cite[Proposition~3.4]{li2022neural} at $p = \frac{1}{2}$ (with $\Delta\tau = 1/N$),
\begin{equation}\label{eq:relu_drift_calc}
\begin{aligned}
\hE[c\,\varphi_s(z_1)\varphi_s(z_2)] &= \rho_\ell^{\alpha\beta} + \nu(\rho_\ell^{\alpha\beta})\,\Delta\tau + \mathcal{O}(\Delta\tau^{3/2}), \\
\hE\!\left[c\,\varphi_s(w^\alpha)\varphi_s(w^\beta) \mid \mathcal{F}_\ell\right] &= \sqrt{\Phi_\ell^{\alpha\alpha}\Phi_\ell^{\beta\beta}}\left(\rho_\ell^{\alpha\beta} + \nu(\rho_\ell^{\alpha\beta})\,\Delta\tau + \mathcal{O}(\Delta\tau^{3/2})\right) \\
&= \Phi_\ell^{\alpha\beta} + b^{\alpha\beta}(\Phi_\ell)\,\Delta\tau + \mathcal{O}(\Delta\tau^{3/2}),
\end{aligned}
\end{equation}
whose remainder divided by $\Delta\tau$ is $\mathcal{O}(\Delta\tau^{1/2}) \approx 0$, yielding~\eqref{eq:drift_relu_like}.

In the smooth case ($s = a\sqrt{N} = a/\sqrt{\Delta\tau}$), Transfer of Taylor's theorem at the origin up to fourth order gives, in conditional internal $L^4$,
\begin{equation}\label{eq:smooth_taylor_L4}
\varphi_s(x) = x + \frac{\varphi''(0)}{2a}\sqrt{\Delta\tau}\,x^2 + \frac{\varphi'''(0)}{6a^2}\Delta\tau\,x^3 + \mathcal{O}_{L^4}(\Delta\tau^{3/2}).
\end{equation}
Taking $g \sim \hNcal(0, 1)$ in~\eqref{eq:smooth_taylor_L4} and using transferred Gaussian moments gives
\begin{equation}\label{eq:smooth_c_expansion}
c^{-1} = 1 + \frac{\Delta\tau}{a^2}\left(\frac{3}{4}\varphi''(0)^2 + \varphi'''(0)\right) + o(\Delta\tau), \qquad c = 1 - \frac{\Delta\tau}{a^2}\left(\frac{3}{4}\varphi''(0)^2 + \varphi'''(0)\right) + o(\Delta\tau),
\end{equation}
where $o(\Delta\tau)$ denotes a remainder $\eta$ satisfying $\eta/\Delta\tau \approx 0$. The odd cubic terms have zero internal expectation, while Transfer of Isserlis' formula gives $\hE[(w^\alpha)^3 w^\beta \mid \mathcal{F}_\ell] = 3\Phi_\ell^{\alpha\alpha}\Phi_\ell^{\alpha\beta}$ and $\hE[(w^\alpha)^2(w^\beta)^2 \mid \mathcal{F}_\ell] = \Phi_\ell^{\alpha\alpha}\Phi_\ell^{\beta\beta} + 2(\Phi_\ell^{\alpha\beta})^2$. Consequently,
\begin{equation}\label{eq:smooth_unnorm_exp}
\begin{aligned}
\hE\!\left[\varphi_s(w^\alpha)\varphi_s(w^\beta) \mid \mathcal{F}_\ell\right] &= \Phi_\ell^{\alpha\beta} + \frac{\varphi''(0)^2}{4a^2}\Delta\tau\left(\Phi_\ell^{\alpha\alpha}\Phi_\ell^{\beta\beta} + 2(\Phi_\ell^{\alpha\beta})^2\right) \\
&\quad + \frac{\varphi'''(0)}{2a^2}\Delta\tau\,\Phi_\ell^{\alpha\beta}\left(\Phi_\ell^{\alpha\alpha} + \Phi_\ell^{\beta\beta}\right) + o(\Delta\tau).
\end{aligned}
\end{equation}
Multiplying~\eqref{eq:smooth_unnorm_exp} by~\eqref{eq:smooth_c_expansion} gives
\begin{equation}\label{eq:smooth_norm_exp}
\hE\!\left[c\,\varphi_s(w^\alpha)\varphi_s(w^\beta) \mid \mathcal{F}_\ell\right] = \Phi_\ell^{\alpha\beta} + b^{\alpha\beta}(\Phi_\ell)\,\Delta\tau + o(\Delta\tau),
\end{equation}
with $b$ exactly as in~\eqref{eq:drift_smooth}.

\medskip\noindent\emph{Covariance calculation.}
For each $\alpha, \beta$, define the centered internal fluctuations
\begin{equation}\label{eq:shaped_R_def}
R_{\ell,i}^{\alpha\beta} := c\,\varphi_s(h_{\ell+1,i}^\alpha)\varphi_s(h_{\ell+1,i}^\beta) - \hE\!\left[c\,\varphi_s(w^\alpha)\varphi_s(w^\beta) \mid \mathcal{F}_\ell\right], \qquad R_\ell^{\alpha\beta} := \frac{1}{\sqrt{N}}\sum_{i=1}^N R_{\ell,i}^{\alpha\beta}.
\end{equation}
Then
\begin{equation}\label{eq:shaped_inc_entry}
\Phi_{\ell+1}^{\alpha\beta} - \Phi_\ell^{\alpha\beta} = b^{\alpha\beta}(\Phi_\ell)\,\Delta\tau + \sqrt{\Delta\tau}\,R_\ell^{\alpha\beta} + o(\Delta\tau),
\end{equation}
where the final term is the conditional-mean remainder from the preceding drift calculation. In both shaping schemes, $c\,\varphi_s(h_{\ell+1,1}^\alpha)\varphi_s(h_{\ell+1,1}^\beta) = h_{\ell+1,1}^\alpha h_{\ell+1,1}^\beta + \mathcal{O}(\sqrt{\Delta\tau})$ in conditional internal $L^2$. After centering, this gives $R_{\ell,1}^{\alpha\beta} = h_{\ell+1,1}^\alpha h_{\ell+1,1}^\beta - \Phi_\ell^{\alpha\beta} + \mathcal{O}(\sqrt{\Delta\tau})$ in conditional internal $L^2$. Conditional on $\mathcal{F}_\ell$, the leading term has the same internal law as the Gaussian-product fluctuation computed in Proposition~\ref{prop:isserlis_gram} and~\eqref{eq:sigma_tensor_def}. Conditional $\hP$-independence across $i = 1, \dots, N$ then gives
\begin{equation}\label{eq:shaped_cov_approx}
\begin{aligned}
\hCov\!\left(R_\ell^{\alpha\beta}, R_\ell^{\gamma\delta} \mid \mathcal{F}_\ell\right) &= \hCov\!\left(R_{\ell,1}^{\alpha\beta}, R_{\ell,1}^{\gamma\delta} \mid \mathcal{F}_\ell\right) \\
&= \hCov\!\left(w^\alpha w^\beta, w^\gamma w^\delta \mid \mathcal{F}_\ell\right) + o(1) \\
&= \Sigma_{\alpha\beta,\gamma\delta}(\Phi_\ell) + o(1),
\end{aligned}
\end{equation}
where $o(1)$ denotes an infinitesimal entry in $\Inf$. Let $R_\ell$ now denote the vector of its $\bar{m}$ distinct upper-triangular entries. Then, in flattened coordinates,
\begin{equation}\label{eq:sigma_N_approx}
\Sigma_N(\Phi_\ell) := \hCov(R_\ell \mid \mathcal{F}_\ell) = \Sigma(\Phi_\ell) + o(1) \approx \Sigma(\Phi_\ell).
\end{equation}
By Transfer, one may choose $\xi_\ell$ with the internal conditional moments in~\eqref{eq:xi_cond_cov} so that
\begin{equation}\label{eq:R_sqrt_factor}
R_\ell = \Sigma_N(\Phi_\ell)^{1/2}\xi_\ell.
\end{equation}
If $\Sigma_N(\Phi_\ell)$ is internally singular, its zero-variance directions do not contribute to $R_\ell$; one can adjoin independent unit-variance coordinates in those directions so that $\hCov(\xi_\ell \mid \mathcal{F}_\ell) = I_{\bar{m}}$ without changing the factorization. By~\eqref{eq:sigma_N_approx} and S-continuity of the positive-semidefinite matrix square root, $\Sigma_N(\Phi_\ell)^{1/2} \approx \Sigma(\Phi_\ell)^{1/2}$, and hence with $\Delta B_\ell = \sqrt{\Delta\tau}\,\xi_\ell$,
\begin{equation}\label{eq:sqrt_swap_remainder}
\sqrt{\Delta\tau}\,R_\ell = \Sigma(\Phi_\ell)^{1/2}\Delta B_\ell + \tilde{o}_\ell(\Delta\tau).
\end{equation}
Combining this identity with~\eqref{eq:relu_drift_calc} or~\eqref{eq:smooth_norm_exp} proves~\eqref{eq:kem_step}. The same transferred Gaussian moment bounds give the infinitesimal conditional third-moment (Lindeberg) rate used in the diffusion shadow below; see~\cite[Appendices~C--D]{li2022neural} for the corresponding uniform estimates on compact sets.
\end{proof}

The expansion~\eqref{eq:kem_step} has the form of one Keisler--Euler--Maruyama step for a diffusion with drift $b$, diffusion covariance $\Sigma$, and diffusion coefficient $\Sigma^{1/2}$. This leads to the following pathwise shadow theorem (compare~\cite{li2022neural}).

\begin{theorem}[Neural Covariance SDE Limit {\cite{li2022neural}}]\label{thm:neural_cov_sde}
Fix standard $m \in \N$ and standard inputs $x^1, \dots, x^m \in \R^{n_0}$. Let $\bar{m} := m(m+1)/2$, let $D, N \in \hNinf$ satisfy $\st(D/N) = \bar{\tau} \in (0, \infty)$, and let $\{\Phi_\ell\}_{\ell=0}^D$ be the internal shaped Gram-matrix Markov chain defined by~\eqref{eq:shaped_gram}. Assume a critical shaping scheme as in Theorem~\ref{thm:kem_increment}, and suppose the standard input Gram matrix $\Phi_0$ is strictly positive definite. Define the internal depth-time interpolation
\begin{equation}\label{eq:shaped_interp}
\Phi_\tau^{(N)} := \Phi_{\min\{\lfloor \tau N \rfloor,\, D\}}, \qquad \tau \ge 0.
\end{equation}
For standard $r > 0$, define the norm-exit stopping times
\begin{equation}\label{eq:norm_exit_times}
\tau_r^{(N)} := \inf\left\{\tau \ge 0 : \norm{\Phi_\tau^{(N)}} > r\right\}, \qquad \tau_r := \inf\left\{\tau \ge 0 : \norm{\Phi_\tau} > r\right\},
\end{equation}
where $\norm{\cdot}$ denotes the Euclidean norm of the $\bar{m}$ upper-triangular entries. Then for every standard radius $r > 0$ at which the stopping map is almost surely continuous for the limiting diffusion (all but at most countably many $r > 0$), the stopped hyperfinite trajectory $\tau \mapsto \Phi_{\tau \wedge \tau_r^{(N)}}^{(N)}$ is $\Loeb$-almost surely limited and S-continuous on $[0, \bar{\tau}]$, and its uniform shadow $\widetilde{\Phi}^{(r)}_\tau := \st(\Phi_{\tau \wedge \tau_r^{(N)}}^{(N)})$ on $(\Omega, \mathcal{L}(\mathcal{A}), \Loeb)$ has the law of the stopped path $\Phi_{\cdot \wedge \tau_r}$ in $C([0, \bar{\tau}], \R^{\bar{m}})$ (lifting directly to pathwise equality
\begin{equation}\label{eq:stopped_shadow_conv}
\st\!\left(\Phi_{\tau \wedge \tau_r^{(N)}}^{(N)}\right) = \Phi_{\tau \wedge \tau_r} \qquad \text{uniformly on } [0, \bar{\tau}]
\end{equation}
under Keisler's Loeb-space representation; see Appendix~A.3), where $\Phi$ is the unique maximal strong solution of
\begin{equation}\label{eq:neural_cov_sde}
d\Phi_\tau = b(\Phi_\tau)\,d\tau + \Sigma(\Phi_\tau)^{1/2}\,dB_\tau, \qquad \Phi_0 \in \SPD(m),
\end{equation}
with drift $b$ and diffusion tensor $\Sigma$ as in Theorem~\ref{thm:kem_increment}, and with $B_\tau$ a standard Brownian motion in $\R^{\bar{m}}$ in the upper-triangular coordinates.
\end{theorem}

\begin{proofsketch}
Flatten $\Phi$ into its $\bar{m}$ upper-triangular entries, as in Section~\ref{sec:multi_cov_sde}. The Keisler--Euler--Maruyama increment expansion in Theorem~\ref{thm:kem_increment}, together with the transferred Gaussian conditional-moment bound at the end of its proof, gives the required drift, covariance, and infinitesimal third-moment rate uniformly on standard compact subsets of $\SPD(m)$. Because these conditional estimates depend on the history only through $\Phi_\ell$, they furnish the Markov-state coefficients required by Theorem~A.4. They match the drift $b$ and diffusion covariance $\Sigma$ in~\eqref{eq:neural_cov_sde}, yielding S-continuity and the stopped shadow up to exit from regular localization domains lying strictly inside $\SPD(m)$ (driven directly on $(\Omega, \mathcal{L}(\mathcal{A}), \Loeb)$ by $B_\tau = \st(\sum_{\ell < \lfloor \tau N \rfloor} \Delta B_\ell)$ prior to boundary contact); the precise sufficient lifting criterion is given in Theorem~A.4. Extending from strict interior domains of $\SPD(m)$ to the norm-ball stopping time $\tau_r$ uses the boundary localization of~\cite[Appendices~A, C, and~D]{li2022neural}.
\end{proofsketch}

An independent readout converts the terminal internal Gram matrix into the output covariance shadow. Let $W_D \in \hR^{1 \times N}$ be an internal readout weight row vector with i.i.d.\ $\hNcal(0, 1)$ entries, $\hP$-independent of the weights used to form $h_D^\alpha$. Define the (scalar) internal readout for input $\alpha$ at hyperfinite depth $D$ by
\begin{equation}\label{eq:shaped_readout}
f^\alpha := \sqrt{\frac{c}{N}} W_D \varphi_s(h_D^\alpha), \qquad \alpha = 1, \dots, m.
\end{equation}
Conditional on $\mathcal{F}_D$, by Transfer the vector $f := (f^1, \dots, f^m) \in \hR^m$ is internally jointly Gaussian with mean zero and covariance
\begin{equation}\label{eq:readout_cond_cov}
\hCov\!\left(f^\alpha, f^\beta \mid \mathcal{F}_D\right) = \frac{c}{N}\ip{\varphi_s(h_D^\alpha)}{\varphi_s(h_D^\beta)} = \Phi_D^{\alpha\beta}.
\end{equation}
Since this internal conditional law depends on $\mathcal{F}_D$ only through $\Phi_D$, it follows that
\begin{equation}\label{eq:readout_given_phi}
f \mid \Phi_D \sim \hNcal(0, \Phi_D).
\end{equation}
Whenever the diffusion~\eqref{eq:neural_cov_sde} does not explode by time $\bar{\tau}$ (so $\Phi_D \in \Lim^{m \times m}$ and $\st(\Phi_D) = \Phi_{\bar{\tau}}$ $\Loeb$-almost surely), $f$ is $\Loeb$-almost surely limited and the shadow readout satisfies $\st(f) \mid \Phi_{\bar{\tau}} \sim \mathcal{N}(0, \Phi_{\bar{\tau}})$, forming a Gaussian scale mixture with random covariance $\Phi_{\bar{\tau}}$~\cite[Theorems~3.2 and~3.9]{li2022neural}.

We close the section by examining stability in the one-input case. When $m = 1$ (one input / one covariance coordinate), the smooth-shaped drift~\eqref{eq:drift_smooth} reduces to a logistic-type scalar SDE:
\begin{equation}\label{eq:scalar_logistic_sde}
d\Phi_\tau = b\,\Phi_\tau(\Phi_\tau - 1)\,d\tau + \sqrt{2}\,\Phi_\tau\,dB_\tau, \qquad b := \frac{\frac{3}{4}\varphi''(0)^2 + \varphi'''(0)}{a^2},
\end{equation}
where $B_\tau$ is a standard one-dimensional Brownian motion.

A useful intuition for the role of $b$ comes from the corresponding deterministic ODE $\partial_\tau \Phi_\tau = b\,\Phi_\tau(\Phi_\tau - 1)$, which solves explicitly to
\begin{equation}\label{eq:scalar_logistic_ode}
\Phi_\tau = \frac{1}{1 - (1 - \Phi_0^{-1})e^{b\tau}}.
\end{equation}
If $b > 0$ and $\Phi_0 > 1$, the ODE explodes at the finite time $\tau_\star = \frac{1}{b}\log\!\bigl(\frac{\Phi_0}{\Phi_0 - 1}\bigr)$, whereas for $\Phi_0 \in (0, 1)$ it is driven toward $0$. If $b < 0$, the drift points toward $1$ from both sides, so every positive ODE trajectory remains bounded in $(0, \infty)$; if $b = 0$, every ODE trajectory is constant. For the SDE, the Brownian term at $\Phi_\tau = 1$ is $\sqrt{2}\,dB_\tau$, so a trajectory that reaches $1$ continues to fluctuate rather than staying there. Nevertheless, the same drift mechanism explains why escape from the positive appreciable galaxy $(0, \infty) = \st(\App_{>0})$ is possible: when $b > 0$, an upward fluctuation above $1$ is reinforced by a positive drift that grows quadratically, whereas for $b \le 0$ the drift does not push large values farther outward. The following result makes this intuition precise.

\begin{proposition}[Finite-Time Explosion Criterion for the Scalar Covariance SDE {\cite[Proposition~3.7]{li2022neural}}]\label{prop:scalar_explosion}
Let $\Phi_\tau$ solve~\eqref{eq:scalar_logistic_sde} with standard $\Phi_0 > 0$, and define the exit time from $(0, \infty)$ (the shadow of the positive appreciable ring $\App_{>0}$) by
\begin{equation}\label{eq:explosion_time_def}
\tau^\star := \sup_{M > 1} \inf\left\{\tau \ge 0 : \Phi_\tau \ge M \text{ or } \Phi_\tau \le M^{-1}\right\}.
\end{equation}
Then $\Loeb(\tau^\star = \infty) = 1$ if and only if $b \le 0$; equivalently, finite-time explosion out of the positive appreciable galaxy $(0, \infty)$ occurs with positive probability if and only if $b > 0$.
\end{proposition}

Thus $b \le 0$ is the exact almost-sure nonexplosion condition keeping the covariance trajectory inside the positive appreciable galaxy; when $b > 0$, finite-time escape from $(0, \infty)$ occurs with positive probability, not necessarily almost surely~\cite[Section~4]{li2022neural}.

This is one reason the smooth-shaped covariance SDE shadow in Theorem~\ref{thm:neural_cov_sde} is naturally formulated in a \emph{local} (stopped) sense: when escape from $\App_{>0}$ is possible, one compares the hyperfinite chain $\Phi_\tau^{(N)}$ and the continuous diffusion $\Phi_\tau$ up to the first exit time from each standard compact subset of $(0, \infty)$ (or of $\SPD(m)$ in the matrix case)~\cite[Definition~3.8 and Theorem~3.9]{li2022neural}.

\section{Exercises}\label{sec:ch5_exercises}

\begin{exercise}[Local Truncation Error]\label{ex:local_trunc_error}
Assume~\eqref{eq:b_lipschitz}. For a positive infinitesimal step size $\eta \approx 0$, show that the one-step error of the internal Euler update~\eqref{eq:euler_scheme} satisfies
\begin{equation}\label{eq:local_trunc_O_eta2}
\eta^{-2}\norm{\ns{X}_\eta - (X_0 + \eta\,b(X_0))} \in \Lim, \qquad \text{i.e.,} \qquad \norm{\ns{X}_\eta - (X_0 + \eta\,b(X_0))} = \mathcal{O}(\eta^2),
\end{equation}
and use Transfer of the discrete Gr\"onwall inequality to deduce~\eqref{eq:euler_uniform_conv} in Theorem~\ref{thm:euler_conv}.
\end{exercise}

\begin{exercise}[Wick Expansion]\label{ex:wick_expansion}
Let $g \sim \hNcal(0, \Phi)$ in $\hR^m$. Use Transfer of Isserlis' theorem to compute $\hE[g_\alpha g_\beta g_\gamma g_\delta]$ and derive Proposition~\ref{prop:isserlis_gram}.
\end{exercise}

\begin{exercise}[Proof of Proposition~\ref{prop:relu_sde_one_sample}]\label{ex:relu_sde_one_sample}
Adapt the one-sample argument of~\cite[\S2.2]{li2022neural} to hyperfinite width $N \in \hNinf$ and depth $D \in \hNinf$ with $\st(D/N) = \bar{\tau}$ and $\Delta\tau := 1/N \approx 0$. In particular, show by Transfer that $\Phi_{\ell+1} = \Phi_\ell\bigl(1 + \sqrt{\Delta\tau}\,R_\ell\bigr)$ for $0 \le \ell \le D - 1$, where $\hE[R_\ell \mid \mathcal{F}_\ell] = 0$ and $\hVar(R_\ell \mid \mathcal{F}_\ell) = 5$. After checking the infinitesimal conditional third-moment rate and applying Doob's maximal inequality~\eqref{eq:doob_maximal}, match the zero drift and diffusion coefficient $\sqrt{5}\,\Phi$ in Proposition~\ref{prop:relu_sde_one_sample} via Theorem~A.4, and verify~\eqref{eq:relu_terminal_log} by summing the second-order nonstandard Taylor expansion of $\ns{\!\log}(1 + \sqrt{\Delta\tau}\,R_\ell)$ over $\ell = 0, \dots, D-1$.
\end{exercise}
